\section{Paper Ideas}
\label{sec:paper-ideas}

Two possible publications have come out of the study so far. Both are
thin in their current form. This section states each idea, the evidence
we have, and what is still missing.

\subsection{Structure-Aware Deflation for Constraint-Space Solvers}
\label{sec:idea-deflation}

\paragraph{Idea.}
Simulators that work in constraint space solve
$\mathbf{A}\boldsymbol{\lambda} = \mathbf{b}$ either iteratively (PCR, CG,
PGS) or directly (Cholesky of $\mathbf{A}$, or of the body-space matrix
$\mathbf{K}$). Deflated PCR (Section~\ref{sec:dpcr}) sits between the two:
a small block is solved exactly and the rest iteratively. The paper would
present DPCR with a $\mathbf{Z}$ that needs no eigenvectors: the
joint/control block (Section~\ref{sec:z-joint-control}) plus $k_c$
contact rows chosen by the AMG strength-of-connection rule
(Section~\ref{sec:z-strength}). The claim is that this hybrid is faster,
or more accurate for the same time budget, than both pure iterative and
pure direct solving on some class of robotics scenes.

\paragraph{Evidence so far.}
On Helhest, the median PCR iteration count drops from $19$ (Jacobi) to
$7$ and $9$ with the joint/control block, and to $5$ and $4$ with $8$
additional strength rows, in the obstacle and terrain scenes
(Figure~\ref{fig:deflation-z-sources}). The GPU time per solve at $W = 1$
drops from about $0.63$ to $0.35$\,ms and from $0.84$ to $0.51$\,ms
(Figure~\ref{fig:deflation-z-sources-time}). $\mathbf{Z}$ costs nothing to
build for the joint/control block and little for the strength rows.

\paragraph{Weaknesses.}
\begin{itemize}
  \item \textbf{The direct alternative wins on Helhest.} The body-space
    Cholesky solve (Section~\ref{sec:body-space-idea}) is $3.9$ to $5.6$
    times faster than the best PCR variant at $W = 1$. A reviewer will ask
    why one should not simply factorize $\mathbf{K}$.
  \item \textbf{The parts are known.} Unit-vector $\mathbf{Z}$ is static
    condensation, the strength rule is the Ruge--St\"uben C/F criterion,
    and AGX and \citet{peiret2019schur} solve joints directly and
    contacts iteratively. Only the combination in a differentiable GPU
    simulator appears to be new.
  \item \textbf{The gain is modest}, about $1.6$ to $1.8$ times at
    $W = 1$, and the iteration counts of Figure~\ref{fig:deflation-z-sources}
    are in double precision while the engine runs in single precision.
\end{itemize}

\paragraph{What it would take.}
Scenes where body space is impossible or expensive: many bodies (long
chains, quadrupeds, piles) so that $(6N_b)^3$ is too large, or rigid rows
with $c_i = 0$. Baselines: Jacobi PCR, a block-Jacobi preconditioner with
the exact joint block (not yet measured; the most obvious competitor),
body-space Cholesky, and a Kamino-style dense solve \citep{tsounis2026kamino}.
The most original element is the observation that both rows of a strongly
coupled pair must be eliminated (Section~\ref{sec:z-strength}); a cheap
selection rule with a bound on the remaining condition number would give
the paper a theoretical contribution.

\paragraph{Venue.}
A simulation workshop at RSS, ICRA or CoRL with the current material plus
two or three scenes; RA-L or IROS with scenes where body space loses.

\clearpage
\subsection{Hybrid Body/Constraint Space for Accurate Adjoint Gradients}
\label{sec:idea-hybrid-adjoint}

\paragraph{Idea.}
A differentiable simulator that eliminates the velocities and solves for
the constraint forces computes the gradient carried to the previous step
as $\mathbf{w}_u = \mathbf{d} - \tilde{\mathbf{M}}^{-1}\mathbf{J}^\top\mathbf{w}$
\eqref{eq:adjoint-pcr}. This is the Woodbury formula for
$\mathbf{K}^{-1}\mathbf{f}$, and its subtraction amplifies the relative
error of $\mathbf{w}$ by $\|\tilde{\mathbf{M}}^{-1}\mathbf{J}^\top\mathbf{w}\| / \|\mathbf{w}_u\|$.
The hybrid eliminates a group of rows $S$ in body space,
$\hat{\mathbf{M}} = \tilde{\mathbf{M}} + \mathbf{J}_S^\top\mathbf{C}_S^{-1}\mathbf{J}_S$,
and keeps the remaining rows $F$ in constraint space,
$\mathbf{A}_F = \mathbf{J}_F\hat{\mathbf{M}}^{-1}\mathbf{J}_F^\top + \mathbf{C}_F$.
It is exact in exact arithmetic and, for $S$ = joint/control rows, its
forward system is the same Schur complement that joint/control deflation
iterates on. The claim: the hybrid restores, in maximal coordinates with
compliant constraints, what minimal-coordinate engines get for free, the
elimination of the joints from the constraint system, and so makes
budget-limited single-precision adjoints accurate.

\paragraph{Evidence so far.}
Offline, single step, single precision, against the double-precision
solution of the same linear system
(\texttt{test\_scripts/deflation/adjoint\_hybrid.py}):
\begin{itemize}
  \item The amplification is about $30$ (median) for the engine's formula
    and $2.5$ with the joint/control rows in body space.
  \item With the engine's PCR budget ($16$ iterations), the share of
    draws with more than $10\%$ gradient error drops from $42\%$ to $2\%$
    (obstacle) and from $25.5\%$ to $2\%$ (terrain), and to $0\%$ with
    one refinement step.
  \item Iterative refinement, the textbook fix of the Woodbury
    instability \citep{yip1986note,hashemi2025instability}, does not
    rescue the engine's formula under the cap ($42\% \to 38.5\%$): the
    error contracts by amplification $\times$ inner error per step, which
    exceeds $1$. With direct solves, one refinement step makes every
    variant accurate to about $10^{-7}$.
\end{itemize}

\paragraph{Scope and prior art.}
The same structure appears in Simple \citep{lelidec2024highly} (forward
mode) and Nimble \citep{werling2021fast}, but both solve the constraint
part directly, probably in double precision, and use minimal coordinates,
in which the joints are not constraint rows. The amplification becomes
harmful only with an inexact constraint solve (an iteration cap, single
precision) and stiff constraints that absorb most of the loss push. The
contribution is therefore the regime where GPU engines are heading:
maximal coordinates, single precision and budgeted matrix-free iterative
solves, where the textbook fix fails.

\paragraph{What it would take.}
\begin{itemize}
  \item Evidence at the gradient level: finite-difference agreement and a
    trajectory optimization with both adjoints inside the engine. This
    needs the hybrid in \texttt{src/} and coordination with the authors
    of the engine.
  \item A second engine: running Simple's derivative formula with an
    iterative single-precision solve would show the effect outside
    ostrich.
  \item A general bound for the amplification in terms of
    $\tilde{\mathbf{M}}$, $\mathbf{J}$ and $\mathbf{C}$, and a rule for
    rows with $c_i = 0$, which cannot enter body space.
\end{itemize}

\paragraph{Venue.}
A differentiable-simulation workshop with the offline results and an
in-engine finite-difference check; a full paper (RA-L, SCA) with
optimization results and a second engine.

\clearpage
\subsection{Selecting the Solver per Scene}
\label{sec:idea-solver-selection}

\paragraph{Idea.}
No single linear solver is expected to be best for every scene. Pure
iterative PCR scales to large systems, a direct body-space solve wins when
there are few bodies and all rows are compliant, and the hybrid of
Sections~\ref{sec:idea-deflation} and~\ref{sec:idea-hybrid-adjoint} covers
the region in between. The idea is to choose the solver automatically,
per scene or robot type, so as to reach the highest accuracy, including
the accuracy of the adjoint gradient, within a given time budget. The
choice could be made by handcrafted rules or learned, for example with a
decision-tree or MLP classifier.

\paragraph{Prior art.}
Learning to choose a linear solver is established in scientific
computing. \citet{sood2019iterative} labels sparse systems by the fastest
of many PETSc solver--preconditioner pairs and trains classifiers on cheap
matrix features; later work uses graph neural networks on the sparsity
pattern and large SuiteSparse benchmarks. These works select for speed
only, one matrix at a time. Not covered so far: a choice that includes
gradient accuracy, batched GPU simulation, and a hybrid option between the
direct and the iterative extreme.

\paragraph{Concerns.}
\begin{itemize}
  \item \textbf{Something to select.} On Helhest the body-space solve
    wins everywhere, so a classifier would always give one answer. The
    idea needs scenes in which each option wins: many bodies, rigid rows
    with $c_i = 0$, contacts between bodies.
  \item \textbf{Cost model versus learning.} The deciding features are
    few and largely analytic: $N_b$, $n$, the size of the joint block,
    $\min_i c_i$, the batch size $W$, the precision and whether gradients
    are needed. The costs follow known formulas ($(6N_b)^3/3$ for a dense
    $\mathbf{K}$, iterations $\times$ cost per iteration for PCR). A cost
    model with a few measured constants may match a classifier. A decision
    tree is defensible because it yields the rule; an MLP for about five
    features is hard to justify.
  \item \textbf{Measuring instead of predicting.} A simulation runs
    millions of steps, so it can run a few steps with every solver at
    startup and keep the best, exactly and for the actual hardware. A
    predictor beats this only if the best choice changes during a
    rollout, for example from few rows in the air to many contacts after
    landing, or to contacts between bodies in a grasp. Switching per step
    or per world conflicts with captured CUDA graphs and with all worlds
    of a batch running the same kernels.
\end{itemize}

\paragraph{What it would take.}
First the regime map: which solver wins on which scene, offline from
rescaled masses and compliances, then on real scenes. Then a constrained
choice over a configuration space, $\{\text{PCR}, \text{hybrid with } k_c
\text{ rows}, \text{body-space Cholesky}\} \times \{\text{iteration cap}\}$,
maximizing accuracy subject to time $\le$ budget. Three selectors to
compare: an analytic cost model, a decision tree fitted to measurements,
and benchmarking at startup. The novel case is switching during a rollout
from cheap per-step features (number of active rows, share of contact
rows, number of contacts between bodies).

\paragraph{Assessment.}
Weak as a standalone paper, because of the prior art and the strong
startup-benchmarking baseline. Strong as the closing section of a paper
on choosing the linear solver in differentiable contact simulation, where
the decision tree turns the regime map into a rule readers can use.
